% Member 4 contribution; insert into the team's CVPR template.
% This is a section fragment, not a standalone document.
% Merge member4_references.bib with the team bibliography.
% Snapshot: 2026-11-16. Formal B/D results remain pending.

\subsection{Photometric augmentation}
We investigate whether modest appearance changes during training improve
transparent-object segmentation under controlled photometric shifts.
Model C uses the same ImageNet-initialized MiT-B0 encoder and newly initialized
SegFormer-style decoder as A~\cite{xie2021segformer}; it does not fine-tune A's
trained checkpoint. With probability 0.5, we uniformly select one operation
from brightness scaling, contrast scaling and gamma correction, with its
parameter sampled uniformly from $[0.8,1.2]$. Operations follow resizing and
precede ImageNet normalization, clip intensities to $[0,1]$, and leave target
and valid-pixel masks unchanged. Contrast uses each image's per-channel mean.
These perturbations change appearance rather than physically simulate refraction.
Checks on 20 fixed training images verified mask invariance, reproducibility
for a fixed seed/epoch/image ID, and agreement between cached and uncached
inputs. Validation receives no random augmentation. C retains A's
BCE-plus-Dice objective with the boundary auxiliary head disabled.

\subsection{Controlled A/C experiment}
The reviewed Trans10K-v1 split~\cite{xie2020transparent} contains 4,959 training,
995 validation and 4,428 test images. Training/validation reductions from the
nominal split follow the documented data audit; official test membership is
unchanged. A/C share initial parameters, seed 42, $512\times512$ inputs,
batch size 8, AdamW with initial learning rate $6\times10^{-5}$, BF16 training,
30 epochs and 18,600 updates. C's epoch-30 checkpoint is selected by foreground
IoU on all validation images, without using test scores for selection.
Inference uses FP32, original-size bilinear logit interpolation, a fixed
sigmoid threshold of 0.5 and no post-processing. Foreground IoU and region F1
pool pixel counts; Boundary IoU~\cite{cheng2021boundary} averages image scores.

Both models evaluate every test image under clean inputs and 12 independent
perturbations: brightness $\{0.6,0.8,1.2,1.4\}$, contrast
$\{0.5,0.75,1.25,1.5\}$, and gamma $\{0.6,0.8,1.2,1.6\}$.
Each condition starts from the original input. The perturbation mean excludes
clean inputs, and the drop equals the clean score minus that mean.

\begin{table}[t]
\centering
\small
\setlength{\tabcolsep}{3pt}
\begin{tabular}{lrrrr}
\hline
 & \multicolumn{2}{c}{Clean} & \multicolumn{2}{c}{Perturbed mean} \\
Metric & A & C & A & C \\
\hline
IoU & 87.4364 & 87.5004 & 86.7838 & 86.9604 \\
F1 & 93.2971 & 93.3336 & 92.9235 & 93.0248 \\
Boundary IoU & 68.0870 & 68.0586 & 66.8749 & 67.0003 \\
\hline
\end{tabular}
\caption{Full-test A/C scores (\%). Each condition covers 4,428 images;
the mean includes 12 perturbations. Formal B/D results are pending.}
\label{tab:member4-ac}
\end{table}

C improves mean perturbed IoU by 0.1766 percentage points (pp), with the
clean-to-perturbed drop decreasing from 0.6526 to 0.5401 pp
(Table~\ref{tab:member4-ac}). IoU is higher in all 12 perturbation conditions.
The mean gain is 0.3990 pp on the 937 hard images (75.9338\% to 76.3329\%),
versus 0.0924 pp on the 3,491 easy images. Clean IoU improves by 0.0640 pp,
but clean Boundary IoU decreases by 0.0283 pp. Thus the observed average
robustness benefit is small and involves a clean-boundary trade-off.

Individual regressions remain: for test/hard/images/765.jpg under gamma 1.6,
A/C IoU is 0.6906/0.0643. This case follows the predefined selection rule of
the three smallest C-minus-A differences under gamma 1.6; it does not
estimate error frequency. These are single-seed results without significance
testing. Training hardware and some auxiliary dependencies differ, as recorded
in the environment comparison. We do not infer universal or statistically
significant gains.

\subsection{Transfer to GSD glass scenes}
We freeze 813 GSD test records~\cite{lin2021glass} from mirror revision
\texttt{f8897574} (full revision in the accompanying provenance file),
preserving source image and mask bytes. A single-channel mask value of 255
denotes glass. Exact and near-duplicate searches against all reviewed
Trans10K splits found no cross-dataset candidates at the recorded thresholds;
the limited search cannot prove complete separation. A 20-image label spot
check and review preceded scoring. We retain all records, including one
duplicate photograph with differing masks (812 distinct decoded images).
The full and supplementary manifests have identical ordered members, sharing
one prediction run rather than constituting independent experiments.

\begin{table}[t]
\centering
\small
\setlength{\tabcolsep}{3pt}
\begin{tabular}{lrrrr}
\hline
Model & IoU & F1 & Boundary IoU & Bg. FPR \\
\hline
A & 68.1273 & 81.0425 & 44.2109 & 11.4190 \\
D & \multicolumn{4}{c}{Pending formal model and evaluation} \\
\hline
\end{tabular}
\caption{GSD scores (\%) on 813 frozen records. Background false-positive
rate (Bg. FPR) is lower-is-better. No GSD training or tuning is performed.}
\label{tab:member4-gsd}
\end{table}

The validation-selected A checkpoint obtains the scores in
Table~\ref{tab:member4-gsd} without GSD training, target-set tuning or threshold
changes. In a selected low-score example, A predicts the central open window
area as glass while the annotation marks the side panes. The gallery contains
20 preselected samples and four additional cases from the predefined
lowest-IoU/Boundary-IoU rule. This experiment measures glass-scene transfer;
it does not replace the separate held-out-category experiment.

\paragraph{Pending comparisons.}
Formal B/D models are still pending. Their delivery and comparability checks
are required before full B/D perturbation evaluation, D's GSD evaluation,
and the four-model and A/D conclusions. The current evidence does not
establish the effect of boundary supervision or its interaction with augmentation.

% Optional contribution statement for the team's required responsibility list:
% Member 4: photometric augmentation verification; formal C training;
% full A/C robustness evaluation; GSD import, audit, manifest freezing and
% A transfer evaluation; reproducible release and report/presentation materials.
